\chapter{Planificación}

\begin{quotation} [Abogado] {Peter Drucker (1909--2005)}
La planificación a la largo plazo no es pensar en decisiones futuras, sino en el futuro de las decisiones presentes.
\end{quotation}

\begin{abstract}
La planificación del proyecto es uno de los mecanismos fundamentales para lograr un objetivo, donde se explica lo que se necesita hacer y cómo debe llevarse a cabo.
\end{abstract}

\section{Resumen temporal del proyecto}

\begin{table*}[htb]
	\centering
	\begin{coolTable}{ll}{2}
{Resumen del proyecto}
	\textbf{Fecha de inicio}&11/11/2024\\
	\textbf{Fecha de fin}&03/06/2025\\
	\textbf{Periodicidad de las revisiones}&2 semanas\\
	\textbf{Carga de trabajo semanal}&14 horas\\
	\textbf{Horas totales previstas}&300 horas\\ % entre 25-30 horas por crédito
	\textbf{Horas finales}&x horas\\
	\end{coolTable}
	\caption{Tabla resumen de tiempos y planificación}
\end{table*}

\section{Planificación inicial}
En el contexto de este proyecto, al hablar de planificación, es necesario destacar que se establecieron fechas claves desde el inicio, que se siguieron con éxito durante el desarrollo del proyecto para garantizar el cumplimiento de los objetivos.

La primera de estas fechas claves fue el comienzo del proyecto, el \textbf{11 de noviembre}, marcando una semana para definir el alcance del proyecto, la planificación inicial y el lenguaje de programación, optando por Node.js y React para el desarrollo del Back-End y Front-End respectivamente.
De la misma manera se establecen iteraciones desde el comienzo del desarrollo del proyecto bisemanales, hasta el, se estima también por tanto decisiones del proyecto cada 3 semanas.

Además, el \textbf{14 de noviembre}, se mantuvo una reunión con la tutora del TFM para presentar el alcance del proyecto, así como definir la planificación del proyecto, estableciendo las fechas de entrega de las distintas partes del proyecto, así como las reuniones de seguimiento del proyecto.

Se establecieron iteraciones bisemanales desde el comienzo del desarrollo del proyecto. Además, se dedicaron los lunes y viernes a sesiones de trabajo conjunto entre los dos desarrolladores del proyecto, priorizando en estos días el desarrollo de la documentación del proyecto.
\begin{table*}[htb]
	\centering
	\begin{coolTable}{ll}{2}
{Resumen de iteraciones}
	\textbf{Iteración 1}&11/11/2024 a 24/11/2024\\
	\textbf{Iteración 2}&25/11/2024 a 08/12/2024\\
    \textbf{Iteración 3}&09/12/2024 a 22/12/2024\\
    \textbf{Iteración 4}&23/12/2024 a 05/01/2025\\
    \textbf{Iteración 5}&06/01/2025 a 19/01/2025\\
    \textbf{Iteración 6}&20/01/2025 a 02/02/2025\\
    \textbf{Iteración 7}&03/02/2025 a 16/02/2025\\
    \textbf{Iteración 8}&17/02/2025 a 02/03/2025\\
    \textbf{Iteración 9}&03/03/2025 a 16/03/2025\\
    \textbf{Iteración 10}&17/03/2025 a 30/03/2025\\
    \textbf{Iteración 11}&31/03/2025 a 13/04/2025\\
    \textbf{Iteración 12}&14/04/2025 a 27/04/2025\\
    \textbf{Iteración 13}&28/04/2025 a 11/05/2025\\
    \textbf{Iteración 14}&12/05/2025 a 25/05/2025\\
	\textbf{Iteración 15}&26/05/2025 a 03/06/2025\\
	\end{coolTable}
	\caption{Planificación temporal de iteraciones con fecha de entrega 3 de junio}
\end{table*}


Durante la primera iteración se planificó el desarrollo de la estructura de la base de datos, la creación de la base de datos.

Para la segunda iteración se planificó la conexión de la base de datos con el Back-End de la aplicación.

Para la tercera iteración se planificó el inicio de sesión por parte del usuario con el sistema de Microsoft.

Durante la cuarta iteración se modificó el desarrollo del inicio sesión y se solicitó la creación de una máquina virtual para el despliegue de la aplicación.
